\section{Pure Axioms and Node Creating Rules}
\label{sec:pure-ext}

A common advantage of proof systems for Hybrid Logics is their ability to establish a very general completeness result, in which it is shown that certain extensions of the base proof system are automatically complete for the corresponding classes of models.  This is in sharp contrast with many proof systems for other Modal Logics, in which simple extensions of a complete proof system might be incomplete.

This property of Hybrid Logics is particularly useful when designing proof systems for specific classes of models. For instance, certain applications may require structural constraints ---such as transitivity or irreflexivity--- on accessibility relations.
This is specially relevant for XPath, as the different axes are usually associated with semantics (e.g., price, sub-category, sub-menu, max-level, etc.) which might need to obey structural restrictions. 
Extensions using pure axioms and node-creating rules address these requirements naturally, while still yielding automatic completeness results.

Extensions for $\Hxpath$ that preserve completeness have been investigated in~\cite{ArecesF21} for Hilbert-style axiom systems. In this section, we adopt a similar approach to develop complete tableau systems for extensions of $\Hxpath$. To this end, we follow the methodology introduced in~\cite{Blackburn00,Blackburn2002BeyondPA,BtC06} for Hybrid Logics. We will consider two kinds of extension: pure axioms, and node creating rules.  These extensions ensure the soundness and completeness of the resulting calculi, but do not generally preserve decidability or complexity bounds. Such results depend on carefully controlling the flow of information within the tableaux and must be established on a case-by-case basis. While some structural properties can be managed using the mechanisms introduced earlier, others may fall outside the scope of such approach.

\subsection{Pure Axioms} 
One of the most standard extensions of hybridized proof systems is via pure axioms. This extension enables us to characterize several interesting classes of models, including some that are not definable by so-called ``orthodox'' node expressions (i.e., expressions containing only propositional symbols). 

\begin{defi}\label{def:pure}
    A node expression is \emph{pure} iff it does not contain propositional symbols.
\end{defi}

\begin{exa}\label{ex:pure}
 Below we list some pure node expressions defining interesting frame conditions:
 \[
 \begin{array}{l@{\quad\quad}l}
    \text{Irreflexivity}  &  i{:}\neg \tup{\child} i \\
    \text{Transitivity} &  i{:}(\tup{\child}\tup{\child}j \ra \tup{\child}j) \\
    \text{Trichotomy} & i{:}\tup{\child} j \lor i{:}j \lor j{:}\tup{\child}i\\
    \text{Data Uniqueness} & \tup{@_i =_\cmp @_j} \to i{:} j .
 \end{array}
 \]
 It is interesting to notice that transitivity is definable without the use of nominals,  but trichotomy and irreflexivity are not (see \cite{Blackburn00}). The latter is also true for data uniqueness.
\end{exa}

It has been shown in~\cite{ArecesF21} that pure expressions correspond to first-order properties when taken as axioms. This makes them suitable to characterize interesting classes of models. Below we formally introduce such a notion.

\begin{defi}
    A \emph{(data) frame} is a tuple $\aframe = \tup{N, \set{R_\pathh}_{\pathh \in \rels}, \set{\approx_{\cmp}}_{\cmp\in\cmps}}$, where each component is as in~\Cref{def:models}.
    We write $\aframe\vDash\varphi$ iff for all $\model$ based on $\aframe$, and $n \in N$, we have $\model,n \vDash \varphi$.
    We say that a node expression~$\varphi$ \emph{defines a class of frames $\fframe$} iff $\aframe\vDash\varphi$ is equivalent to  $\aframe\in\fframe$. 
\end{defi}

Our tableaux receive as input node expressions that are prefixed by some $i{:}$, so in order to handle pure axioms we need to make sure they are also prefixed. Following~\cite{Blackburn00,Blackburn2002BeyondPA}, in the following lemma we show that this is not a real restriction.

\begin{lem}\label{lemma:defines}
 If a pure node expression $\varphi$ defines a class of frames $\fframe$ and $i$ a nominal not in $\varphi$, then, $i{:}\varphi$ also defines~$\fframe$.
\end{lem}
\begin{proof}
    The result follows immediately from the requisite that $i$ does not appear in $\varphi$, as in this case $i$ does not affect the meaning of $\varphi$.
\end{proof}

We are now in position to state a general completeness result. The idea is to incorporate (prefixed) pure expressions to check satisfiability of a node expression in a particular class of models. In other words, if we look for a model of a given node expression, pure axioms will guarantee that this model will belong to the corresponding class of models. Given a set of pure expressions $\axioms$ and $\psi\in\axioms$, we need to instantiate nominals appearing in $\psi$, with nominals appearing in the branch of the tableaux in which~$\psi$ is being incorporated. For instance, for a branch $\Theta$ with $k{:}\varphi\in\Theta$, we incorporate irreflexivity as formulated in~\Cref{ex:pure} with the instance $k{:}\neg\tup{\child}k$. This kind of remark applies also to data comparisons. 

\begin{thm}[Strong Completeness for Pure Axioms]\label{th:strong}
    \label{th:pure-completeness}
    Let $\axioms$ be a set of pure node expressions prefixed by some $i{:}$.
    Extend the tableau calculus from~\Cref{sec:tableaux} with the rule (Pure) which adds, at any stage, any $\psi\in\axioms$ instantiated with nominals already appearing in the branch. The obtained tableau is complete with respect to the class of models based on frames defined by $\bigwedge_{\psi\in\axioms} \psi$.
\end{thm}

\begin{proof}
 The argument for this proof is the same as the one in~\cite[Th.~7.2]{Blackburn00} and \cite[Prop.~3.2]{Blackburn2002BeyondPA} for tableau calculi for hybrid logics, and~\cite[Th.~3.23]{ArecesF21} for axiomatizations of hybrid XPath.  Intuitively, the rule (Pure) ensures that pure axioms in $\axioms$ are fully instantiated in a suitable way, in an open and saturated branch $\Theta$. As a consequence, the extracted model $\model^\Theta$ globally satisfies all relevant instances in $\axioms$. This is sufficient to ensure that the frame of $\model^\Theta$ satisfies the structural properties associated with $\axioms$. In 
 other words, the extracted model $\model^\Theta$, satisfying the root formula in the branch, will belong to the required model class. 
\end{proof}

By choosing a suitable set of pure axioms it is possible to characterize interesting XPath features. 
We introduce some of them in the following example.

\begin{exa}\label{ex:extensions}
As our first example of the use of pure actions, we begin by considering the  \emph{inverse} relation, a distinctive feature in XPath fragments which allows for backwards navigation. To this end, we extend the set of relation symbols in the signature to $\rels'=\rel\cup\set{\child^- \mid \child \in \rels}$. 
%Suppose that for each $\child\in\rels$, we have a symbol $\child^-\in \rels$.  
Then, we define
\[
     \axioms[Inverse] = \{i{:}(j\to[\child]\tup{\child^-}j),i{:}(j\to[\child^-]\tup{\child}j) \mid \child \in \rels\}.
\]
By incorporating $\axioms[Inverse]$, our tableaux yields a complete calculus for $\Hxpath$ with inverse relations (over the new signature). More precisely, pure axioms in $\axioms[Inverse]$ ensure that, in any model considered, $R_\textsf{a} = (R_{\textsf{a}^-})^-$. Hence, $\textsf{a}^-$ can be used to refer to the inverse of 
$\textsf{a}$. 

For a second example, fix now $\rels''=\rels'\cup\set{\sib}$, and define
\[
     \axioms[Sibling] = \axioms[Inverse] \cup \{i{:}(\tup{\sib}j\to \tup{\child^-}\tup{\child}j), i{:}(\neg j \land \tup{\child^-}\tup{\child}j \to \tup{\sib}j), i{:}(j{:}\neg\tup{\sib}j) \mid \child \in \rels\}.
\]
By incorporating $\axioms[Sibling]$ we get, additionally, completeness for $\Hxpath$ with $R_{\sib}$ being the \emph{sibling} relation. The first node expression states that a sibling is one node found going backwards and then forwards in the model. The second node expression states the converse of the first expression, in case $j$ is not the current node. Finally, the last expression expresses that a node is not its own sibling. 

As a final example, define 
\[
    \axioms[\sib\text{-}Irreflexivity] = \axioms[Sibling] \cup \{ i{:}(\tup{\sib} j \to (\tup{@_i \neq_\cmp @_j})) \}.  
\]
$\axioms[\sib\text{-}Irreflexivity]$ gives us completeness with respect to the class of models in which siblings have different data values for the comparison criteria $\cmp$. This could be useful, for example, if when the siblings in the data tree represent categories with different prize ranges (e.g., economy, business, first-class).
\end{exa}

We stress that the Strong Completeness Theorem (\Cref{th:strong}) guarantees that the resulting calculi are complete for the intended model classes.  All the extensions discussed in~\Cref{ex:extensions} correspond to common features found in various versions of XPath. It is therefore desirable to have tableau calculi that are not only sound but also complete for these extensions. As it turns out, in all these cases, completeness is automatically preserved due to the specific form of the axioms being added. This being said, adding such axioms does not imply complexity bounds are preserved.  
%In fact, complexity guarantees often fail to hold for these cases. 
For instance, the satisfiability problem for Hybrid Logic $\mathcal{H}(:)$ extended with the inverse operator is already known to be \textsc{ExpTime}-complete~\cite{Areces2000c}. Since $\mathcal{H}(:)$ with inverses is a proper fragment of $\Hxpath$ extended with inverses, this immediately yields an \textsc{ExpTime} lower bound for the satisfiability problem of $\Hxpath$ with inverses.

\subsection{Node Creating Rules} We now turn our attention to introducing so-called \emph{node creation rules} into our tableaux. To precisely characterize the expressivity provided by these new rules, we need to extend the basic language $\Hxpath$ to a strong version with quantification over nominals, as is done for Hybrid Logic~\cite{BlackburnS95}. It is important to clarify that the strong hybrid language for XPath introduced in~\Cref{def:strong} is used solely to characterize the expressive power of the node-creating rules defined below. It is also worth noting that the satisfiability problem for the strong hybrid language for XPath is undecidable~\cite{arec:hybr05b}. Therefore, node-creating rules must be used with caution when aiming for a tableau calculus with good computational properties. 

\begin{defi}[Strong Hybrid Language for XPath]\label{def:strong}
    \label{def:stronghl}
    The \emph{strong hybrid language} for XPath is obtained by allowing expressions of the form $\exists x.\varphi$ (where $x\in\noms$, and $\varphi$ is a node expression) as node expressions.
    The semantics of these new node expressions over models $\model = \tup{N, \set{R_\pathh}_{\pathh \in \rels}, \set{\approx_{\cmp}}_{\cmp\in\cmps}, \valuation, \naming}$ is given as:
    \[
        \begin{array}{lcl}
            \model,n\vDash\exists x.\varphi & \text{\it iff } & \text{there is $ \naming'$ s.t. } \model',n\vDash \varphi,
        \end{array}
    \]
    where $\model'= \tup{N, \set{R_\pathh}_{\pathh \in \rels}, \set{\approx_{\cmp}}_{\cmp\in\cmps}, \valuation, \naming'}$ for  $\naming'$ defined exactly as $\naming$ except, perhaps, for $x$. The expression $\forall x.\varphi$ is defined as $\neg\exists.\neg\varphi$.  We will use $\exists x,y,\ldots$ and $\forall x,y,\ldots$ as shortcuts for $\exists x.\exists y\ldots$ and $\forall x.\forall y\ldots$, respectively.
\end{defi}

The strong hybrid language helps highlight a key limitation of pure axioms and motivates the need for node-creating rules. Pure axioms correspond to universally quantified expressions in the strong hybrid language (see~\cite[Prop.~3.2]{Blackburn2002BeyondPA} for details). However, certain natural frame properties cannot be captured using pure axioms alone.
A well-known example is the Church-Rosser class of frames, or the property of right-directness. The former property can be expressed as 
\[
    \forall i,j,k.\exists x. (i{:}\tup{\child} j \land  i{:}\tup{\child}x \to j{:}\tup{\child}x \land k{:}\tup{\child}x),
\]
while the latter can be expressed as
\[
    \forall i,j.\exists k. (i{:}\tup{\child} k \land j{:}\tup{\child} k).
\]
Both expressions involve a $\forall\exists$ quantification pattern, which goes beyond the purely universal pattern allowed by pure axioms. Node-creating rules provide a natural way to handle such $\forall\exists$ properties: the premises correspond to the universally quantified parts, while the conclusions create new nodes serving as witnesses for the existential quantifiers.

For example,
\[
\forall i,j.\exists k. \tup{@_i\neq_\cmp @_k}\land\tup{@_j\neq_\cmp @_k}
\]
can be casted into a node generating rule as
\[
\begin{array}{c}
{\infer[\text{with $k$ a new nominal.}]
{
\deduce[]
{\phantom{|} {\tup{@_i\neq_\cmp @_k}}}    
{\phantom{|}{\tup{@_j\neq_\cmp @_k}}}}
{\text{$i,j$ in the branch}
}
}
\end{array}
\]
Incorporating this rule into our tableaux ensures that for every pair of nodes $i$ and $j$ in a branch, there is a node $k$ whose data value differs from both.

Let us now move to the general form of these rules and how to represent them. Let $\forall x_1,\ldots,x_\ell.\exists y_1,\ldots y_{\ell'}.\varphi$ be a node expression from the strong hybrid language, such that $\varphi$ is quantifier free. For every such a node expression we define the rule  
\[
\begin{array}{c}
{\infer[\text{with $k_1,\ldots,k_{\ell'}$ new nominals.}]
{
\phantom{|} {\varphi[x_1 \leftarrow i_1,\ldots,x_\ell \leftarrow i_\ell, y_1\leftarrow k_1\ldots y_{\ell'}\leftarrow k_{\ell'}]}}
{\text{$i_1,\ldots,i_{\ell}$ in the branch}}
}
\end{array}
\]

Once more, we can prove a Strong Completeness Theorem for Node Creating Rules.

\begin{thm}[Strong Completeness for Node Creating Rules]
    Let $S$ be a set of node expressions from the strong hybrid language of the form $\forall x_1,\ldots,x_\ell.\exists y_1,\ldots y_{\ell'}.\varphi$, where $\varphi$ is quantifier free.
    Extend the tableau calculus from~\Cref{sec:tableaux} with the associated rules as defined above. The obtained tableaux is complete with respect to the class of models whose frames are defined by $\bigwedge_{\psi\in S} \psi$.
\end{thm}

\begin{proof}
The argument is as in the proof of~\Cref{th:strong}.
\end{proof}